\chapter*{Introduction}
\addcontentsline{toc}{chapter}{Introduction}
\VerdierPhysicalPageAnchor{0018}
\VerdierSemanticAnchor{intro}

\VerdierSemanticAnchor{intro-opening}
In this work we propose a formalism for hyperhomology.

\VerdierSemanticAnchor{intro-1}
\noindent\textbf{1.}\quad
Let $\mathcal A$ and $\mathcal B$ be two abelian categories, and let
$F:\mathcal A\to\mathcal B$ be an additive functor. When $\mathcal A$ has
enough injective objects, the construction of the right derived functors of
$F$ is now classical [1]: every object $X$ of $\mathcal A$ admits an injective
resolution $X\xrightarrow{\varepsilon}I^{\bullet}(X)$ unique ``up to
homotopy''. The cohomology objects of the complex $F(I^{\bullet}(X))$ are
therefore determined by $X$ up to canonical isomorphism and depend
functorially on $X$. The resulting functors are the right derived functors of
$F$.

\VerdierSemanticAnchor{intro-1-homotopy-quasi}
Thus, to each object $X$ of $\mathcal A$ one functorially associates a complex
of objects of $\mathcal B$ defined ``up to homotopy'', whose cohomology is then
taken. Most often, however, the usual constructions of homological algebra do
not provide complexes defined up to homotopy, but complexes defined
\emph{up to quasi-isomorphism} (a morphism of complexes is called a
quasi-isomorphism if it induces an isomorphism on cohomology objects). To
illustrate this point, let us give two examples.

\VerdierSemanticAnchor{intro-1-f-acyclic}
Let us first note that, in order to study the derived functors of $F$, one
should consider resolutions $X\xrightarrow{\varepsilon'}\mathcal R(X)$ more
general than injective resolutions: resolutions by objects acyclic for $F$
(in sheaf theory, resolutions by flabby or soft sheaves [23]; in module
theory, resolutions by flat modules, $\ldots$). These resolutions are no
longer necessarily unique up to homotopy. But given an $F$-acyclic resolution
$X\xrightarrow{\varepsilon'}\mathcal R(X)$ and an injective resolution
$X\xrightarrow{\varepsilon}I^{\bullet}(X)$, there exists a unique (up to
homotopy) morphism of resolutions, i.e. a morphism of complexes
$\rho:\mathcal R(X)\to I^{\bullet}(X)$ such that the following\ %
